\section{Discussion}
\label{sec:discussion}

This section critically evaluates the three discussed architectural approaches for implementing the Visitor pattern. The comparison focuses on runtime performance, structural design, adherence to SOLID principles, and ease of maintenance.

\subsection{Classic Visitor Pattern}

\textbf{Advantages:}
\begin{itemize}
    \item \textbf{Extensibility of Operations:} It facilitates the seamless addition of new operations (Visitors) without modifying the existing codebase of the element hierarchy, strictly adhering to the Open/Closed Principle (OCP) in terms of behaviors \cite{gof1994}.
    \item \textbf{Cross-Hierarchy Processing:} The \texttt{Visitor} interface can declare \texttt{visit} methods for concrete elements that do not share the same abstract base class, allowing operations to span multiple disparate class hierarchies.
\end{itemize}

\textbf{Disadvantages:}
\begin{itemize}
    \item \textbf{Runtime Overhead:} The double-dispatch mechanism necessitates two virtual function calls (vtable lookups) per element. This indirection inherently slows down execution compared to a standard single dispatch.
    \item \textbf{Cache Inefficiency:} Classic elements are typically managed via pointers and allocated dynamically on the heap. This scattered memory layout causes frequent CPU cache misses, further degrading runtime performance.
    \item \textbf{Intrusive Architecture:} Implementing this pattern requires injecting the \texttt{accept} method into every element class. This makes it an unviable solution if the developer lacks access to the original source code (e.g., when dealing with third-party libraries).
    \item \textbf{Type-Addition Bottleneck:} The design is fundamentally unsuited for systems where new element types are frequently added. Introducing a new concrete element requires modifying the base \texttt{Visitor} interface, which subsequently forces a recompilation of all existing visitor implementations, violating the OCP for types.
    \item \textbf{Circular Dependency and Code Duplication:} The pattern introduces a rigid circular dependency between the element header files and the visitor header files. Furthermore, it often forces developers to write boilerplate code, repeating similar \texttt{visit} logic for different elements.
\end{itemize}


\subsection{Modern C++ Approach (\texttt{std::variant} and \texttt{std::visit})}

\textbf{Advantages:}
\begin{itemize}
    \item \textbf{Observed Runtime Advantage in the Benchmark:} By storing objects by value, \texttt{std::variant} ensures contiguous memory allocation (cache-friendly). Furthermore, it utilizes compile-time type-based visitation over a closed set of alternatives rather than traditional vtables, reducing dispatch overhead in common scenarios \cite{cpprefvariant,cpprefvisit}.
    \item \textbf{Non-Intrusive Design:} The elements remain entirely independent. There is no need to define a base class or implement an \texttt{accept} method, allowing this approach to work seamlessly with third-party or legacy classes.
    \item \textbf{Flexible Subsetting:} Developers can define a variant using a specific subset of derived classes rather than being forced to depend on the entire hierarchy.
    \item \textbf{Code Deduplication:} By leveraging generic lambdas (\texttt{auto}) or templated \texttt{operator()} methods inside the visitor functor, developers can handle multiple types implicitly, significantly reducing repetitive code (adhering to the DRY principle).
\end{itemize}

\textbf{Disadvantages:}
\begin{itemize}
    \item \textbf{Violation of the Dependency Inversion Principle (DIP):} Unlike traditional OOP where high-level modules depend on abstractions (e.g., \texttt{Shape*}), \texttt{std::variant} requires hardcoding the exact list of concrete types (e.g., \texttt{std::variant<Circle, Rectangle>}). This couples the design directly to concrete implementations.
    \item \textbf{Type-Addition Bottleneck (OCP Violation):} Similar to the Classic Visitor, adding a new type requires modifying the \texttt{std::variant} type alias and updating the functor to handle the new type, which triggers extensive recompilation.
\end{itemize}


\subsection{The Acyclic Visitor Pattern}

\textbf{Advantages:}
\begin{itemize}
    \item \textbf{Decoupled Architecture:} It entirely eliminates the circular dependency problem found in the Classic Visitor by utilizing a degenerate base interface \cite{martin1997}.
    \item \textbf{Safe and Independent Evolution:} Acyclic Visitor reduces the coupling between the central visitor interface and individual element types, allowing new element-specific visitor interfaces to be introduced without modifying a monolithic Visitor interface. This prevents cascading compilation errors when extending parts of the system.
    \item \textbf{Selective Implementation:} Concrete visitors are only required to inherit and implement the specific visitor interfaces for the elements they actually intend to process, avoiding bloated ``catch-all'' interfaces.
\end{itemize}

\textbf{Disadvantages:}
\begin{itemize}
    \item \textbf{Severe Runtime Overhead:} The core dispatch mechanism relies on cross-casting via \texttt{dynamic\_cast}. Run-Time Type Information (RTTI) lookups are notoriously slow in C++, making this pattern unsuitable for performance-critical applications (e.g., game engines or high-frequency trading) \cite{cpprefrtti}.
    \item \textbf{Retained Intrusiveness:} Despite fixing the dependency cycle, the Acyclic Visitor still requires element classes to implement the \texttt{accept} method, maintaining an intrusive footprint on the data structures.
\end{itemize}
